\section{同一实验室时刻怎样读取加速杆两端的钟}
比较加速杆两端的钟，需要先指定读取事件。加速坐标的同时面和实验室的同时面选取不同的事件对，钟面读数也随之改变。以下令杆的固有长度保持不变，比较实验室在同一时刻读取的两端固有时。

设杆后端固有加速度为$\alpha$，固有长度$L$，记$s=1+\alpha L/c^2>1$。本讲加速坐标给后端$ct=(c^2/\alpha)\sinh(\alpha\tau/c)$；前端则为$ct=s(c^2/\alpha)\sinh(\alpha\tau/c)$。在两端都取实验室$t=0$读数为零后，前端固有时还要乘局部钟率因子$s$。

\begin{example}同时拍摄两只加速钟的坐标读数
理想记录装置已经消除光传播延迟，得到实验室同一时刻$t$的后端钟读数$T_A$和前端钟读数$T_B$。两钟按上述保持固有长度的方案从静止开始持续加速。求两读数，并判断哪只钟读数更大。这里讨论坐标同时读数，不讨论未经延迟修正的照片。
\ExpansionDiagram{fig:dep-m2-clock-q}{实验室水平同时线选取两条加速世界线上的不同加速坐标时刻}{
\draw[->](0,0)--(4,0)node[right]{$x$};\draw[->](0,0)--(0,3.8)node[above]{$ct$};\draw[thick,domain=0:1.8,samples=70]plot({cosh(\x)-1},{sinh(\x)});\draw[thick,domain=0:1.6,samples=70]plot({1.5*cosh(\x)-1},{1.5*sinh(\x)});\draw[dashed](0,2)--(3.3,2)node[right]{同一$t$};\node[below]at(0,0){$A$};\node[below]at(.5,0){$B$};}
\end{example}
\begin{solution}
图\ref{fig:dep-m2-clock-q}要求先分别反解世界线时间。记$s=1+\alpha L/c^2>1$。后端固有时就是其加速坐标时间，前端固有时为局部钟率因子$s$乘其坐标时间，所以
\[
 T_A=\frac c\alpha\operatorname{arsinh}\frac{\alpha t}{c},\qquad
 T_B=\frac{sc}{\alpha}\operatorname{arsinh}\frac{\alpha t}{sc}.
\]
不能把$T_B$直接写为$sT_A$，因为那样比较的是相同加速坐标时间，不是图中的实验室水平切片。
\ExpansionDiagram{fig:dep-m2-clock-a}{初始读数相同，随后前端钟在实验室相同时间内积累更多固有时}{
\draw[->](0,0)--(5,0)node[right]{$\alpha t/c$};\draw[->](0,0)--(0,3)node[above]{$\alpha T/c$};\draw[thick,domain=0:4.5,samples=90]plot(\x,{ln(\x+sqrt(1+\x^2))});\draw[dashed,domain=0:4.5,samples=90]plot(\x,{1.5*ln(\x/1.5+sqrt(1+(\x/1.5)^2))});\node[right]at(3.5,2.7){$T_B$};\node[right]at(3.5,1.7){$T_A$};}
两式求导，
\[
 \frac{\dd T_A}{\dd t}
 =\frac c\alpha\frac{\alpha/c}{\sqrt{1+(\alpha t/c)^2}}
 =\frac1{\sqrt{1+(\alpha t/c)^2}},
\]
\[
 \frac{\dd T_B}{\dd t}
 =\frac{sc}\alpha\frac{\alpha/(sc)}{\sqrt{1+(\alpha t/(sc))^2}}
 =\frac1{\sqrt{1+(\alpha t/(sc))^2}}.
\]
因$s>1$，对$t>0$有$\dot T_B>\dot T_A$，结合初始两读数为零，得到$T_B>T_A$。图\ref{fig:dep-m2-clock-a}画$s=1.5$的曲线。初期展开可得$T_B-T_A\simeq\alpha^2t^3(1-s^{-2})/(6c^2)$，差值不是从启动瞬间就线性增长。

物理上，保持固有长度要求前端加速度较小。在同一实验室时刻，前端速率较低，钟慢效应较弱，所以固有时积累更多。这一解释依赖本题的具体加速安排；若两端使用相同固有加速度，世界线和所比较距离都改变，不能照搬这些读数。
\end{solution}
